\documentclass[11pt,a4paper]{article}

\usepackage[a4paper,top=1.9cm,bottom=1.9cm,left=2.1cm,right=2.1cm]{geometry}
\usepackage[T1]{fontspec}
\usepackage{lmodern}
\usepackage{microtype}
\usepackage{titlesec}
\usepackage{enumitem}
\usepackage[hidelinks]{hyperref}

\titleformat{\section}{\normalsize\bfseries}{}{0pt}{}
\titlespacing{\section}{0pt}{5pt}{2pt}
\setlist[enumerate]{nosep,leftmargin=1.5em,topsep=2pt}
\setlength{\parindent}{0pt}
\setlength{\parskip}{3pt}
\renewcommand{\refname}{\normalsize References}

\begin{document}

\begin{center}
{\large\bfseries Vidruma: From \textit{Varga} to Vector Space}\\[2pt]
{\normalsize Sense-Disambiguated Distributional Semantics for Sanskrit\\ and the Computational Evaluation of a Traditional Lexical Ontology}\\[3pt]
{\small Sai Sreemayee Sistla (2025114013), Vinayak Sharma (2025114002)}
\end{center}

\vspace{-4pt}

\section{Introduction and Motivation}

    Ideas echoing distributional semantics, and that words get meaning from their company and discourse-context have existed for long in the Indian grammatical tradition \cite{Harris} \cite{Firth} \cite{Bhartrhari} \cite{Anandavardhana} \cite{Matilal}.
    Thesauri like the \textit{Amarako\'sa} have described and categorised words based on ontological principles into twenty-five \textit{varga}s.

    Word-Embedding models have been developed for Sanskrit but report a high number of OOV tokens, upto 55\% \cite{Sandhan}.
    We wish to investigate this and analyse the effect of ID-based lemmatisation on OOV rates.
    To this end, we propose to construct Word-Embedding models, based on lemmatised Sanskrit corpora, and test them against existing benchmarks for Sanskrit Word-Embeddings.
    
    We also wish to compare how modern computational methods can add to the categorisation attempts of traditional lexicographers, and analyse clustering patterns that are encoded in word vectors.
    Lastly, we plan to qualitatively analyse how lemmatisation improves word-sense disambiguation and detecting synonymy.

    Thus, we plan to not only test modern approaches on traditional parameters via existing benchmarks, but also ask how patterns that exist in higher-dimensional vector spaces can complement age-old Sanskrit semantic studies.

\section{Research Questions}

    \begin{enumerate}
        \item[\textbf{RQ1}] How does integrating lemmatisation into word-embeddings layers (such as GloVe and Word2Vec) \cite{Jurafsky} impact semantic categorisation metrics on low-resource Sanskrit benchmarks?
        \item[\textbf{RQ2}] Do unsupervised clusters in the induced space recover \textit{Amarako\'sa varga}s, and what do systematic divergences indicate - unrecorded groupings, cross-\textit{varga} affinities, or limits of the distributional signal across concrete and abstract semantic domains?
        \item[\textbf{RQ3}] How does lemmatisation informed from pragmatic context help resolve word-sense ambiguities between words that have the same lexical realisations? How much of the reported OOV rate for the benchmark is attributable to inflection rather than lexical sparsity?
    \end{enumerate}

\section{Methods}

    We attempt to construct word-embeddings over DCS lemma IDs.
    Preprocessing conditions (surface, transliteration-normalised, string-lemmatised, ID-lemmatised) are taken into account, isolated and measured towards their improvements to the metrics and contribution to the OOV rate.
    Clusters are induced by $k$-means and agglomerative linkage and scored by purity against \textit{varga}s, with divergent cases inspected qualitatively.

    Limitations of the cosine scale across embedding spaces will be taken into account. Wherever needed, additional rank statistics (Spearman $\rho$, pairwise AUC, etc), following critiques of similarity-based evaluation will be used \cite{Schnabel} \cite{Faruqui}.

\section{Datasets and Benchmarking}

    \begin{itemize}
        \item The Digital Corpus of Sanskrit \cite{Hellwig} supplies lemmatised, sense-annotated text ($\approx$3.8M tokens in prior work \cite{Sandhan}).
        \item The \textit{Amarako\'sa} is taken in the Hyderabad digitisation \cite{Nair}: 4{,}056 synsets across 25 \textit{varga}s.
        \item EvalSan \cite{Sandhan} provides four intrinsic tasks with published baselines: relatedness (4{,}551 test pairs balanced over same-synset, same-\textit{varga} and unrelated), synonym detection (3{,}034 items), and semantic and syntactic categorisation (15 and 45 categories).
        We benchmark against the reported results of \cite{Sandhan}.
    \end{itemize}

\begin{thebibliography}{99}\small
\setlength{\itemsep}{0pt}

\bibitem{Harris} Harris, Z. 1954. Distributional structure. \textit{Word} 10(2--3).
\bibitem{Firth} Firth, J.R. 1957. A synopsis of linguistic theory, 1930--1955. \textit{Studies in Linguistic Analysis}.
\bibitem{Sandhan} Sandhan, J., Adideva, O., Komal, D., Behera, L., Goyal, P. 2023. Evaluating neural word embeddings for Sanskrit. \textit{arXiv:2104.00270}; EvalSan toolkit.
\bibitem{Nair} Nair, S., Kulkarni, A. 2010. The knowledge structure in \textit{Amarako\'sa}. \textit{Sanskrit Computational Linguistics}, LNCS 5406.
\bibitem{Hellwig} Hellwig, O. 2010--2021. \textit{The Digital Corpus of Sanskrit (DCS)}.
\bibitem{Bhartrhari} Bhart\d{r}hari. \textit{V\=akyapad\=iya}. Trans.\ K.A.S. Iyer, Deccan College, 1965--1974.
\bibitem{Matilal} Matilal, B.K. 1990. \textit{The Word and the World: India's Contribution to the Study of Language}. Oxford University Press.
\bibitem{Anandavardhana} \=Anandavardhana. \textit{Dhvany\=aloka}. Trans.\ D.H.H. Ingalls, J.M. Masson, M.V. Patwardhan, Harvard Oriental Series 49, 1990.
\bibitem{Schnabel} Schnabel, T., Labutov, I., Mimno, D., Joachims, T. 2015. Evaluation methods for unsupervised word embeddings. \textit{EMNLP}.
\bibitem{Faruqui} Faruqui, M., Tsvetkov, Y., Rastogi, P., Dyer, C. 2016. Problems with evaluation of word embeddings using word similarity tasks. \textit{RepEval@ACL}.
\bibitem{Jurafsky} Jurafsky, D., Martin, J.H. \textit{Speech and Language Processing}, 3rd ed.\ draft, Ch.~6: Vector Semantics and Embeddings.

\end{thebibliography}

\end{document}
